% !TeX root = ../../main.tex

Der strukturelle Aufbau des Systems ist in Abbildung~\ref{fig:softwareDesign} dargestellt und wird im
nachfolgenden Abschnitt detailliert beschrieben.
% !TeX root = ../main.tex

\begin{figure}[ht]
	\centering
	\begin{tikzpicture}
		\node[
			fill=TerminalBackground,
			draw=TerminalFrame,
			rounded corners,
			inner sep=10pt
		] {
			\begin{tikzpicture}[
				scale=0.8,
				transform shape,
				node distance=1cm,
				every node/.style={font=\small},
				startstop/.style={
					rectangle,
					rounded corners,
					draw,
					minimum width=3cm,
					minimum height=1cm,
					align=center
				},
				process/.style={
					rectangle,
					draw,
					minimum width=3cm,
					minimum height=1cm,
					align=center
				},
				decision/.style={
					diamond,
					draw,
					align=center,
					aspect=2
				},
				storage/.style={
					rectangle,
					draw,
					dashed,
					align=center,
					aspect=2
				},
				arrow/.style={
					->,
					>=stealth
				},
				implements/.style={
					dashed,
					-{Triangle[open]}
				},
			]

				\node (start) [startstop] {main.py\\Programmstart};
				\node (runner) [process, below=of start] {experimentRunner.py};
				\node (analysis) [process, right=2cmof runner] {analysis.py};
				\node (experiment) [process, below=of runner] {experimentFunctions.py};
				\node (config) [process, left=2cm of runner] {config.py};
				\node (database) [process, below=of config] {database.py};
				\node (retry) [process, below=of experiment] {retry.py};
				\node (service) [process, below=of database] {service.py};
				\node (log) [storage, below=of analysis] {Logs};

				\draw [arrow, shorten <=5pt, shorten >=5pt] (start) -- (runner);
				\draw [arrow, shorten <=5pt, shorten >=5pt] (start) -| (analysis);
				\draw [arrow, shorten <=5pt, shorten >=5pt] (runner) -- (experiment);
				\draw [implements, shorten <=5pt, shorten >=5pt] (config) -- (experiment);
				\draw [arrow, shorten <=5pt, shorten >=5pt] (experiment.south) to[bend right=20] (retry.north);
				\draw [arrow, shorten <=5pt, shorten >=5pt] (retry.north) to[bend right=20] (experiment.south);
				\draw [arrow, shorten <=5pt, shorten >=5pt] (retry.west) to[bend left=20] (service.east);
				\draw [arrow, shorten <=5pt, shorten >=5pt] (service.east) to[bend left=20] (retry.west);
				\draw [arrow, shorten <=5pt, shorten >=5pt] (service.north) to[bend left=20] (database.south);
				\draw [arrow, shorten <=5pt, shorten >=5pt] (database.south) to[bend left=20] (service.north);
				\draw [implements, shorten <=5pt, shorten >=5pt] (database) -- (experiment);
				\draw [implements, shorten <=5pt, shorten >=5pt] (log) -- (analysis);
				\draw [arrow, shorten <=5pt, shorten >=5pt] (experiment) -- (log);

			\end{tikzpicture}};
	\end{tikzpicture}

	\caption{Aufbau der Systemstruktur}
	\label{fig:softwareDesign}
\end{figure}

Der Programmablauf beginnt mit der Ausführung der \texttt{main}-Funktion.
Von dort aus werden die Funktionen aus dem Modul \texttt{experimentRunner} und die
\texttt{start\_analysis}-Funktion nacheinander aufgerufen.
Die Funktionen des Moduls \texttt{experimentRunner} dienen dabei ausschließlich als namensgebende Funktionen, aus denen
im weiteren Verlauf die verschiedenen Fehlerklassen abgeleitet werden.
Dadurch wird klar zugeordnet, welche Fehlerklassen erzeugt und durch die Experimente getestet werden.
Die anschließende Analyse der Experimente erfolgt über die im Modul \texttt{analysis} enthaltenen Funktionen,
welche die erzeugten Ergebnisse weiterverarbeiten, auswerten und strukturiert darstellen.

Die in \texttt{experimentRunner} implementierten Funktionen rufen die \texttt{start()}-Funktion im Modul
\texttt{experimentFunctions} auf.
Hier sind neben den Experimenten auch die Funktionen zum \emph{Warm-up} der Datenbank und zum Schreiben der Logs
implementiert.
Zusätzlich beziehen sich die Experimente die Parameter aus dem \texttt{config} Modul.
In diesem Modul werden die Clients gestartet, die den Retry ausführen und für jeden Client wird eine neue
Datenbankverbindung aufgebaut.
Die Clients starten daraufhin den Retry mit einer bestimmten Strategie.

Im \texttt{retry}-Modul befinden sich die Implementierungen der einzelnen Retry-Strategien.
In jeder Strategie wird als Erstes die Transaktion ausgeführt.
Dafür ruft eine Funktion den Service auf, der zu der Fehlerklasse passt, den die Funktion aus
\texttt{experimentRunner} am Anfang definiert hat.
In dem Service steht die Implementierung der Transaktion, die den Fehler erzeugen soll.
Diese wird dann von der Datenbank ausgeführt und diese gibt dann entweder den SQLSTATE-Code \texttt{00000} oder den
Fehlerklassen spezifischen Fehler-Code.
Diese werden dann in \texttt{retry} ausgewertet und gespeichert.
Daraufhin setzen die einzelnen Retry-Strategien ein und führen einen Retry an den gegebenen Parametern aus.
Sobald eine Strategie keinen Retry mehr durchführen kann, sei es wegen Überschreitung des maximalen Retry-Counts oder
weil die Transaktion erfolgreich war, werden die Ergebnisse im JSONL-Format an die Funktion in
\texttt{experimentFunctions} zurückgegeben, wo diese dann in die Logdatei geschrieben werden.

Separat ist das \texttt{analysis} Modul.
Diesem wird ein Dateipfad zu den Logs übergeben und wertet diese aus.
Die Ergebnisse gibt es daraufhin strukturiert aus.
